%% =============================================================
%% Section 4: Extensions
%% =============================================================

\section{Extensions}
\label{sec:extensions}

\subsection{Intermediate Communication Regimes}
\label{sec:intermediate}

The binary comparison between forward guidance and opacity abstracts from a range
of communication practices used by actual central banks. Two intermediate regimes
are of particular interest. The \emph{Swiss National Bank} (SNB) publishes
projections of inflation and output conditional on the policy rate remaining constant
throughout the forecast horizon. The \emph{Bank of England} (BOE) publishes
projections conditional on the market-implied path of future interest rates---the
path inferred from financial market instruments at the time of the announcement. In
both cases, the central bank publishes rich economic forecasts, but it does not
directly announce its intended rate path.

\paragraph{The SNB regime.}
Let $\bar{i}$ be the reference rate at which the SNB conditions its forecast; in
practice this is the current period-1 rate $i_1$.\footnote{The SNB's conditioning
assumption has evolved over time, but the essential feature is that the rate is held
fixed at a known, publicly observed level---not the bank's intended future path.}
The announced period-2 forecasts are:
\begin{align}
  \hat{y}_2^{\mathrm{SNB}}  &\equiv \mathbb{E}_1\bigl[y_2 \mid i_2 = i_1\bigr]
    = -\tilde{\sigma}\,i_1, \label{eq:y2hat_SNB}\\
  \hat{\pi}_2^{\mathrm{SNB}} &\equiv \mathbb{E}_1\bigl[\pi_2 \mid i_2 = i_1\bigr]
    = -\kappa\tilde{\sigma}\,i_1 + \varepsilon_c, \label{eq:pi2hat_SNB}
\end{align}
where we have used $\mathbb{E}_1[\varepsilon_{d,2}] = 0$ and the reduced-form
expressions \eqref{eq:IS2_reduced}--\eqref{eq:NKPC2}.

\paragraph{The BOE regime.}
Under opacity, the private sector expects $i_2 = \bar{\phi}_c\,\varepsilon_c$
(equation~\eqref{eq:E_i2_OP}). The BOE conditions its forecast on this
market-implied path:
\begin{align}
  \hat{y}_2^{\mathrm{BOE}}  &\equiv \mathbb{E}_1\bigl[y_2 \mid i_2 =
    \bar{\phi}_c\,\varepsilon_c\bigr]
    = -\tilde{\sigma}\,\bar{\phi}_c\,\varepsilon_c, \label{eq:y2hat_BOE}\\
  \hat{\pi}_2^{\mathrm{BOE}} &\equiv \mathbb{E}_1\bigl[\pi_2 \mid i_2 =
    \bar{\phi}_c\,\varepsilon_c\bigr]
    = \bigl(1 - \kappa\tilde{\sigma}\bar{\phi}_c\bigr)\varepsilon_c.
    \label{eq:pi2hat_BOE}
\end{align}
The conditioning assumption is endogenous to market expectations: private agents
form $\mathbb{E}_1[i_2]$, the BOE conditions on this, and equilibrium requires
consistency between the two. Below we verify that the opacity expectation is the
unique self-consistent fixed point.

\begin{proposition}[Intermediate Regimes Are Welfare-Equivalent to Opacity]
\label{prop:intermediate}
Under either the SNB or the BOE communication regime, the private sector's
equilibrium expectation of the period-2 rate satisfies
$\mathbb{E}_1[i_2] = \bar{\phi}_c\,\varepsilon_c$, the period-1 allocation is
identical to that under opacity, and welfare satisfies
\begin{equation}
  \mathcal{L}^{\mathrm{SNB}} = \mathcal{L}^{\mathrm{BOE}} = \mathcal{L}^{\mathrm{OP}}.
  \label{eq:intermediate_welfare}
\end{equation}
\end{proposition}

\begin{proof}
\textit{SNB regime.} Inspect~\eqref{eq:y2hat_SNB}--\eqref{eq:pi2hat_SNB}. Both
announced forecasts depend only on $i_1$, $\tilde{\sigma}$, $\kappa$, and
$\varepsilon_c$. Of these, $\tilde{\sigma}$ and $\kappa$ are common knowledge,
$i_1$ is publicly observed in period 1, and $\varepsilon_c$ is publicly observed
by assumption (Section~\ref{sec:timing}). The announcements are therefore redundant
given the agents' information set: they contain no information about $\lambda$ beyond
what is already known. Agents continue to use the prior to form
$\mathbb{E}_1^{\mathrm{SNB}}[i_2] = \bar{\phi}_c\,\varepsilon_c$, as under opacity.

\textit{BOE regime.} Suppose agents form expectation
$\mathbb{E}_1[i_2] = m\,\varepsilon_c$ for some coefficient $m$ to be determined.
The BOE conditions on $i_2 = m\,\varepsilon_c$ and announces
\eqref{eq:y2hat_BOE}--\eqref{eq:pi2hat_BOE} with $\bar{\phi}_c$ replaced by $m$.
These announcements depend only on $m$, $\tilde{\sigma}$, $\kappa$, and
$\varepsilon_c$---none of which reveal $\lambda$. Agents learn nothing about
$\lambda$ and update $m$ only through prior beliefs: $m = \bar{\phi}_c$.
The fixed point $m = \bar{\phi}_c$ is therefore unique, confirming
$\mathbb{E}_1^{\mathrm{BOE}}[i_2] = \bar{\phi}_c\,\varepsilon_c$.

Since both intermediate regimes induce the same period-1 expectation as opacity, the
period-1 exchange rate, IS curve, and NKPC (equations
\eqref{eq:q1_regime}--\eqref{eq:NKPC1_r}) are identical. Period-2 policy is
unaffected by period-1 announcements. Welfare is therefore identical to opacity.
\qed
\end{proof}

\paragraph{The key mechanism.}
The result rests on a single structural observation: the central bank's private
information $\lambda$ governs its \emph{reaction function}---how much it raises the
rate in response to cost-push shocks---not the structural equations that describe
how the economy responds to a given rate. The IS and Phillips curves (with slopes
$\tilde{\sigma}$ and $\kappa$) are common knowledge. Conditional forecasts, by
construction, fix the rate path and apply these common-knowledge structural
equations; they therefore mechanically strip out the type-dependent information.
Only direct communication about the intended rate path encodes $\lambda$.

\paragraph{Welfare hierarchy.}
Combining Propositions~\ref{prop:welfare} and~\ref{prop:intermediate}:
\begin{equation}
  \mathcal{L}^{\mathrm{FG}}
  \;\leq\;
  \mathcal{L}^{\mathrm{SNB}}
  \;=\;
  \mathcal{L}^{\mathrm{BOE}}
  \;=\;
  \mathcal{L}^{\mathrm{OP}},
  \label{eq:welfare_hierarchy}
\end{equation}
with strict inequality on the left whenever $p \in (0,1)$, $\lambda_H \neq
\lambda_L$, and $\sigma_c^2 > 0$. The two intermediate regimes, while reflecting
genuine and substantive differences in central bank practice, are
welfare-equivalent to full opacity in this model. The implication is stark: what
matters for welfare is not the \emph{volume} of information published, but whether
the publication encodes information about the central bank's \emph{intended rate
path} and, through it, its type.

\paragraph{Remark: private cost-push shocks.}
The equivalence result depends on $\varepsilon_c$ being publicly observed. If,
instead, $\varepsilon_c$ were private information of the central bank, SNB-style and
BOE-style announcements would reveal $\varepsilon_c$ to private agents (via
\eqref{eq:pi2hat_SNB} and \eqref{eq:pi2hat_BOE}), yielding a strict welfare gain
relative to full opacity. In this richer setting, the intermediate regimes would
occupy a genuinely intermediate position in the welfare hierarchy, transmitting
information about the central bank's macroeconomic assessment without disclosing its
reaction function. We treat this as a natural extension.

%% -----------------------------------------------------------
\subsection{Welfare-Improving Scenario Communication}
\label{sec:scenarios}

Proposition~\ref{prop:intermediate} identifies a sharp boundary: conditional
economic forecasts are welfare-neutral relative to opacity, while explicit rate-path
announcements strictly improve welfare. Between these extremes lies
\emph{scenario communication}---publishing the intended rate path under multiple
contingencies, imperfectly or with deliberate ambiguity. We now parametrise these
intermediate regimes and characterise the welfare gain.

\paragraph{Definition.}
A \emph{scenario} is a pair $\bigl(\varepsilon_c^{(k)},\, \hat{\imath}_2^{(k)}\bigr)$
specifying the central bank's intended period-2 rate conditional on the cost-push
shock taking value $\varepsilon_c^{(k)}$. A \emph{scenario set} is a menu of $K$
such pairs $\mathcal{S} = \bigl\{\bigl(\varepsilon_c^{(k)},\hat{\imath}_2^{(k)}\bigr)
\bigr\}_{k=1}^{K}$. The SNB and BOE regimes are special cases with $\hat{\imath}_2^{(k)}$
fixed independently of $\lambda$; a fully contingent scenario has
$\hat{\imath}_2^{(k)} = \phi_c(\lambda)\varepsilon_c^{(k)}$.

\paragraph{Scenarios and type revelation.}
Since $i_2^*(\lambda) = \phi_d\varepsilon_d + \phi_c(\lambda)\varepsilon_c$, a
scenario $(\varepsilon_c^{(k)},\hat{\imath}_2^{(k)})$ with
$\hat{\imath}_2^{(k)} = \phi_c(\lambda)\varepsilon_c^{(k)}$ at any
$\varepsilon_c^{(k)} \neq 0$ fully reveals $\lambda$ (since $\phi_c$ is strictly
monotone), and is therefore informationally equivalent to full forward guidance.
Non-contingent scenarios (SNB, BOE) fix $\hat{\imath}_2^{(k)}$ independently of
$\lambda$ and hence convey zero type information.

\paragraph{The $\theta$-parametrisation.}
To bridge these extremes we characterise the informativeness of an arbitrary scenario
set $\mathcal{S}$ by the \emph{Blackwell informativeness ratio}:
\begin{equation}
  \theta(\mathcal{S})
  \;\equiv\;
  \frac{\mathrm{Var}_a\bigl[\,\mathbb{E}[\phi_c(\lambda)\mid a]\,\bigr]}
       {\mathrm{Var}_\lambda\bigl[\phi_c(\lambda)\bigr]},
  \label{eq:theta_def}
\end{equation}
where $a$ denotes the signal (announcement) received by private agents and
expectations are taken with respect to the equilibrium signal distribution. By the
law of total variance, $\theta(\mathcal{S}) \in [0,1]$:
\begin{itemize}
  \item $\theta = 0$: no type information is transmitted (opacity, SNB, BOE);
  \item $\theta = 1$: type is fully revealed (full forward guidance);
  \item $\theta \in (0,1)$: partial transmission.
\end{itemize}

\begin{proposition}[Welfare Gain Is Linear in Informativeness]
\label{prop:scenario}
For any scenario set $\mathcal{S}$ with Blackwell informativeness $\theta(\mathcal{S})$:
\begin{equation}
  \mathcal{L}^{\mathrm{OP}} - \mathcal{L}^{\mathcal{S}}
  \;=\;
  \theta(\mathcal{S})\,
  \bigl(\mathcal{L}^{\mathrm{OP}} - \mathcal{L}^{\mathrm{FG}}\bigr).
  \label{eq:scenario_gain}
\end{equation}
The welfare gain from scenario communication is exactly $\theta(\mathcal{S})$ times
the maximum (full forward guidance) gain.
\end{proposition}

\begin{proof}
From the period-1 welfare decomposition~\eqref{eq:L1_gap_expected}, the expected
loss under any communication regime $r$ satisfies
\begin{equation}
  \mathbb{E}[\mathcal{L}_1^r]
  \;=\;
  \mathbb{E}[\mathcal{L}_1^{\mathrm{FG}}]
  + c \cdot \mathbb{E}_a\bigl[\mathrm{Var}[\phi_c(\lambda)\mid a]\bigr]\cdot\sigma_c^2,
  \label{eq:L1_general}
\end{equation}
for some constant $c > 0$ that depends only on the structural parameters. (Under FG
the posterior variance is zero; under opacity it is $\mathrm{Var}_\lambda[\phi_c(\lambda)]$.)
By the law of total variance:
\begin{equation}
  \mathbb{E}_a\bigl[\mathrm{Var}[\phi_c(\lambda)\mid a]\bigr]
  = \mathrm{Var}_\lambda[\phi_c(\lambda)]
    - \mathrm{Var}_a\bigl[\mathbb{E}[\phi_c(\lambda)\mid a]\bigr]
  = (1-\theta)\,\mathrm{Var}_\lambda[\phi_c(\lambda)].
  \label{eq:total_var}
\end{equation}
Substituting~\eqref{eq:total_var} into~\eqref{eq:L1_general}:
\begin{equation}
  \mathbb{E}[\mathcal{L}_1^{\mathcal{S}}]
  = \mathbb{E}[\mathcal{L}_1^{\mathrm{FG}}]
    + c\,(1-\theta)\,\mathrm{Var}_\lambda[\phi_c(\lambda)]\,\sigma_c^2.
\end{equation}
At $\theta = 0$ (opacity): $\mathbb{E}[\mathcal{L}_1^{\mathrm{OP}}] =
\mathbb{E}[\mathcal{L}_1^{\mathrm{FG}}] + c\,\mathrm{Var}_\lambda[\phi_c(\lambda)]\,\sigma_c^2$.
Subtracting:
\begin{equation*}
  \mathbb{E}[\mathcal{L}_1^{\mathrm{OP}}] - \mathbb{E}[\mathcal{L}_1^{\mathcal{S}}]
  = \theta\,c\,\mathrm{Var}_\lambda[\phi_c(\lambda)]\,\sigma_c^2
  = \theta\,\bigl(\mathbb{E}[\mathcal{L}_1^{\mathrm{OP}}]
                  - \mathbb{E}[\mathcal{L}_1^{\mathrm{FG}}]\bigr).
\end{equation*}
Period-2 losses are identical across regimes (the central bank re-optimises in
period~2 regardless of what it announced in period~1), so the equality carries
through to total expected losses. \qed
\end{proof}

\paragraph{Canonical implementation: the $\theta$-disclosure rule.}
A convenient mechanism that achieves any target $\theta \in [0,1]$ is the
\emph{$\theta$-disclosure rule}: the central bank commits to revealing its true
type (announcing $\phi_c(\lambda)\varepsilon_c$) with probability $\theta$, and
making no type-revealing statement (announcing $\bar{\phi}_c\varepsilon_c$) with
probability $1-\theta$. In the binary type model, where
$\phi_c(\lambda_H)\varepsilon_c \neq \phi_c(\lambda_L)\varepsilon_c$ whenever
$\varepsilon_c \neq 0$, agents perfectly infer the regime on the equilibrium path
(a revealing announcement is distinguishable from the prior-mean announcement), so
the rule is credibly incentive-compatible.

Under this rule, the economy is in the FG regime with probability $\theta$ and in
the opacity regime with probability $1-\theta$. Expected total loss:
\begin{equation}
  \mathcal{L}^\theta
  = \theta\,\mathcal{L}^{\mathrm{FG}} + (1-\theta)\,\mathcal{L}^{\mathrm{OP}},
  \label{eq:theta_mixture}
\end{equation}
which immediately yields~\eqref{eq:scenario_gain}. Equation~\eqref{eq:theta_mixture}
makes clear that the welfare hierarchy is a line segment connecting the opacity and
FG outcomes as $\theta$ moves from 0 to 1.

\paragraph{Diminishing returns to signal precision.}
While the welfare gain is linear in informativeness $\theta$, it exhibits
\emph{diminishing returns to signal precision} when the scenario is implemented via
a noisy rate announcement. Suppose the central bank announces
$s = \phi_c(\lambda)\varepsilon_c + \xi$ where $\xi$ is mean-zero noise with
variance $1/\tau$ (precision $\tau \geq 0$). With a Gaussian prior
$\phi_c(\lambda) \sim \mathcal{N}(\bar{\phi}_c, \sigma_\phi^2)$, the posterior
variance is $\sigma_\phi^2/(1 + \tau\sigma_\phi^2\varepsilon_c^2)$, and the
resulting informativeness is:
\begin{equation}
  \theta(\tau) = \frac{\tau\sigma_\phi^2\varepsilon_c^2}{1 + \tau\sigma_\phi^2\varepsilon_c^2},
  \label{eq:theta_tau}
\end{equation}
which is increasing and strictly concave in $\tau$. Combining with
Proposition~\ref{prop:scenario}, the welfare gain as a function of precision is:
\begin{equation}
  \mathcal{L}^{\mathrm{OP}} - \mathcal{L}^\tau
  = \frac{\tau\sigma_\phi^2\varepsilon_c^2}{1 + \tau\sigma_\phi^2\varepsilon_c^2}
    \cdot \bigl(\mathcal{L}^{\mathrm{OP}} - \mathcal{L}^{\mathrm{FG}}\bigr).
  \label{eq:welfare_tau}
\end{equation}
Marginal gains from increasing precision $\tau$ are positive throughout but
converging to zero as $\tau \to \infty$. The full FG gain is recovered only in the
limit $\tau \to \infty$. This concavity implies that \emph{once a central bank has
committed to some scenario communication, the incremental value of further
precision---tighter rate ranges, more scenarios---declines}.

\paragraph{Taxonomy.}
Proposition~\ref{prop:scenario} implies the following welfare-ordered taxonomy:

\begin{enumerate}[label=\arabic*.]
  \item \textbf{Forward guidance} ($\theta = 1$): maximum gain; full type
        revelation via the rate-path announcement.
  \item \textbf{Contingent scenarios} ($\theta \in (0,1)$): intermediate gain
        $= \theta \times$ FG gain; welfare improvement is proportional to
        Blackwell informativeness.
  \item \textbf{Market-conditioned or constant-rate forecasts} ($\theta = 0$):
        zero gain; non-contingent announcements cannot encode the reaction function.
  \item \textbf{Opacity} ($\theta = 0$): benchmark; same welfare as SNB/BOE.
\end{enumerate}

Optimal scenario design (absent commitment frictions) requires $\theta = 1$:
Proposition~\ref{prop:scenario} shows no interior maximum. If the central bank
faces a cost $C(\theta)$ of transparency (legal constraints, political exposure,
or equilibrium credibility bounds), the optimal $\theta^*$ solves:
\begin{equation}
  \theta^* = \arg\max_{\theta\in[0,1]}\;
  \theta\,\bigl(\mathcal{L}^{\mathrm{OP}} - \mathcal{L}^{\mathrm{FG}}\bigr)
  - C(\theta),
  \label{eq:opt_theta}
\end{equation}
yielding an interior solution $\theta^* \in (0,1)$ whenever $C'(0) < \mathcal{L}^{\mathrm{OP}}
- \mathcal{L}^{\mathrm{FG}}$ and $C$ is convex. In this case the marginal value of
transparency equals the marginal cost of disclosure.

%% -----------------------------------------------------------
\subsection{Continuous Type Distribution}
\label{sec:extensions:continuous}

The binary type assumption ($\lambda \in \{\lambda_L, \lambda_H\}$) can be relaxed to
a continuous distribution $\lambda \sim F$ on $[\lambda_L, \lambda_H]$. Under
forward guidance, the announcement $\hat{\imath}_2(\lambda) = \phi_c(\lambda)\varepsilon_c$
fully reveals $\lambda$ for any realisation of $\varepsilon_c \neq 0$, since $\phi_c$
is strictly monotone. The welfare comparison in Proposition~\ref{prop:welfare}
extends directly: the welfare gap between regimes equals
$\mathrm{Var}_\lambda[\phi_c(\lambda)] \cdot \sigma_c^2$, now computed over the
continuous distribution. Richer distributional assumptions (e.g.\ fat tails for the
hawkish type) can affect the quantitative magnitude of the gain.

%% -----------------------------------------------------------
\subsection{Strategic Announcements}
\label{sec:strategic}

In the baseline, the announcement is truthful by assumption. An important extension
considers the case in which the central bank is \emph{strategic}: it chooses what to
announce to influence period-1 private-sector behaviour, subject to the constraint
that the announcement must be credible in a Bayesian equilibrium. Under full
commitment, the truthful announcement remains an equilibrium (a separating equilibrium
in the signalling game), since the announcement perfectly identifies the type. Under
partial commitment, pooling equilibria may arise in which the central bank shades its
announcement toward the prior mean, reducing the informational content of forward
guidance.

%% -----------------------------------------------------------
\subsection{Risk Premia and UIP Deviations}
\label{sec:riskpremia}

The baseline model assumes that UIP holds exactly. A standard extension introduces a
time-varying risk premium $\rho_t$ in the UIP condition:
\begin{equation}
  q_t = \mathbb{E}_t[q_{t+1}] - (i_t - i^*) + \rho_t.
\end{equation}
If $\rho_t$ is correlated with the central bank's type (e.g.\ uncertainty about the
reaction function generates a term premium), then forward guidance has an additional
channel: by resolving type uncertainty, it reduces the risk premium, appreciates the
exchange rate, and further improves period-1 allocations. This channel is
quantitatively relevant for economies with high exchange rate pass-through to import
prices.
